% =========================================================
%  ppt/14-playtest-presentasi.tex
%  Materi 14 — Evaluasi Prototipe, Playtesting, dan Presentasi
%  Bab XVI buku ajar
% =========================================================
\documentclass[aspectratio=169,10pt]{beamer}
% =========================================================
%  ppt/config/packages.tex
%  Paket bersama untuk seluruh presentasi kuliah
%  Pengantar Pemrograman Game
% =========================================================

\usepackage[T1]{fontenc}
\usepackage[utf8]{inputenc}
\usepackage[indonesian]{babel}
\usepackage{lmodern}
\usepackage{textcomp}
\usepackage{amsmath}
\usepackage{amssymb}
\usepackage{graphicx}
\usepackage{booktabs}
\usepackage{tabularx}
\usepackage{array}
\usepackage{multicol}
\usepackage{listings}
\usepackage{upquote}
\usepackage{adjustbox}
\usepackage{tcolorbox}
\usepackage{tikz}
\usetikzlibrary{
  arrows.meta,
  positioning,
  shapes.geometric,
  shapes.misc,
  calc,
  fit,
  backgrounds,
  matrix,
  chains,
  decorations.pathreplacing,
  patterns
}

\tcbuselibrary{skins}

% Jangan memuat geometry: ukuran halaman diatur oleh beamer.

% =========================================================
%  ppt/config/theme.tex
%  Satu tema Beamer untuk keempat belas presentasi
%  Pengantar Pemrograman Game
% =========================================================

% --------------------------------------------------------- warna
\definecolor{ubNavy}{HTML}{1F3864}
\definecolor{ubAccent}{HTML}{2E75B6}
\definecolor{ubHead}{HTML}{D9E2F3}
\definecolor{ubSoft}{HTML}{F2F6FC}
\definecolor{ubLine}{HTML}{BFBFBF}
\definecolor{ubText}{HTML}{1F2937}
\definecolor{ubMuted}{HTML}{5A6678}
\definecolor{ubGreen}{HTML}{1E7B45}
\definecolor{ubGreenBg}{HTML}{F1F8F2}
\definecolor{ubRed}{HTML}{B3261E}
\definecolor{ubRedBg}{HTML}{FDF2F1}
\definecolor{ubWarn}{HTML}{A35A05}
\definecolor{ubWarnBg}{HTML}{FFF8EC}
\definecolor{ubGray}{HTML}{F6F6F6}
\definecolor{ubCodeBg}{HTML}{10233F}
\definecolor{ubCodeText}{HTML}{E8EEF7}
\definecolor{ubCodeKey}{HTML}{8FBCF0}
\definecolor{ubCodeCom}{HTML}{8FB37A}
\definecolor{ubCodeStr}{HTML}{D9A46B}
\definecolor{ubCodeNum}{HTML}{7E99BE}

% Nama warna yang dipakai gambar TikZ buku ajar.
\definecolor{ubPrimary}{HTML}{1F3864}
\definecolor{ubAlt}{HTML}{F2F6FC}

% --------------------------------------------------------- dasar
\usetheme{default}
\usefonttheme{professionalfonts}
\setbeamertemplate{navigation symbols}{}
\setbeamerfont{frametitle}{size=\Large,series=\bfseries}
\setbeamerfont{framesubtitle}{size=\small}
\setbeamerfont{footline}{size=\tiny}
\setbeamerfont{block title}{size=\small,series=\bfseries}
\setbeamerfont{block body}{size=\small}
\setbeamerfont{itemize/enumerate body}{size=\small}
\setbeamerfont{itemize/enumerate subbody}{size=\footnotesize}

\setbeamersize{text margin left=0.60cm,text margin right=0.60cm}
\uchyph=0

\setbeamercolor{normal text}{fg=ubText,bg=white}
\setbeamercolor{structure}{fg=ubAccent}
\setbeamercolor{frametitle}{fg=ubNavy,bg=white}
\setbeamercolor{title}{fg=ubNavy}
\setbeamercolor{footline}{fg=ubMuted,bg=white}
\setbeamercolor{itemize item}{fg=ubAccent}
\setbeamercolor{itemize subitem}{fg=ubAccent}
\setbeamercolor{block title}{fg=white,bg=ubNavy}
\setbeamercolor{block body}{fg=ubText,bg=ubSoft}
\setbeamercolor{block title alerted}{fg=white,bg=ubWarn}
\setbeamercolor{block body alerted}{fg=ubText,bg=ubWarnBg}
\setbeamercolor{block title example}{fg=white,bg=ubGreen}
\setbeamercolor{block body example}{fg=ubText,bg=ubGreenBg}
\setbeamercolor{caption name}{fg=ubMuted}

\setbeamertemplate{blocks}[rounded][shadow=false]
\setbeamertemplate{itemize item}{\raisebox{0.10em}{\color{ubAccent}\scriptsize$\blacksquare$}}
\setbeamertemplate{itemize subitem}{\raisebox{0.06em}{\color{ubAccent}\tiny$\blacktriangleright$}}
\setbeamertemplate{enumerate item}{%
  \makebox[1.3em][l]{\color{ubAccent}\insertenumlabel.}}

\addtobeamertemplate{itemize/enumerate body begin}{%
  \setlength{\itemsep}{0.42em}\setlength{\parskip}{0pt}}{}

% --------------------------------------------------------- judul frame
\setbeamertemplate{frametitle}[default][left]

% --------------------------------------------------------- kaki
\newcommand{\materino}{00}
\newcommand{\materifoot}{Pengantar Pemrograman Game \textperiodcentered\ Materi \materino}

\setbeamertemplate{footline}{%
  \begin{beamercolorbox}[wd=\paperwidth,ht=0.015cm,dp=0cm]{footline}%
    \color{ubHead}\rule{\paperwidth}{0.4pt}%
  \end{beamercolorbox}%
  \begin{beamercolorbox}[wd=\paperwidth,ht=0.40cm,dp=0.14cm,%
                         leftskip=0.60cm,rightskip=0.60cm]{footline}%
    \usebeamerfont{footline}\color{ubMuted}\materifoot\hfill\insertframenumber%
  \end{beamercolorbox}%
}

% --------------------------------------------------------- makro buku ajar
\providecommand{\code}[1]{\texttt{\detokenize{#1}}}
\providecommand{\istilah}[1]{\emph{#1}}
\providecommand{\CPMK}[1]{CPMK\,#1}
\providecommand{\CPL}[1]{CPL\,#1}
\providecommand{\SUBCPMK}[1]{Sub-CPMK\,#1}

% --------------------------------------------------------- gaya kode
\lstdefinelanguage{CSharp}{
  morekeywords={abstract,as,base,bool,break,byte,case,catch,char,class,const,
    continue,decimal,default,delegate,do,double,else,enum,event,explicit,
    extern,false,finally,fixed,float,for,foreach,goto,if,implicit,in,int,
    interface,internal,is,lock,long,namespace,new,null,object,operator,out,
    override,params,private,protected,public,readonly,ref,return,sbyte,
    sealed,short,sizeof,stackalloc,static,string,struct,switch,this,throw,
    true,try,typeof,uint,ulong,unchecked,unsafe,ushort,using,virtual,void,
    volatile,while,var,get,set},
  sensitive=true,
  morecomment=[l]{//},
  morecomment=[s]{/*}{*/},
  morestring=[b]",
}

\lstdefinestyle{kode}{
  basicstyle=\ttfamily\footnotesize\color{ubCodeText},
  keywordstyle=\color{ubCodeKey}\bfseries,
  commentstyle=\color{ubCodeCom}\itshape,
  stringstyle=\color{ubCodeStr},
  numberstyle=\tiny\color{ubCodeNum},
  numbers=left,
  numbersep=7pt,
  backgroundcolor=\color{ubCodeBg},
  frame=none,
  showstringspaces=false,
  breaklines=false,
  tabsize=4,
  columns=fullflexible,
  keepspaces=true,
  xleftmargin=12pt,
  xrightmargin=2pt,
  aboveskip=0pt,
  belowskip=0pt,
  upquote=true,
}
\lstset{style=kode}

\newcommand{\kodejelas}[1]{%
  \vspace{0.16cm}%
  {\footnotesize\color{ubMuted}#1\par}%
}
\newcommand{\sumber}[1]{%
  \vspace{0.08cm}%
  {\scriptsize\color{ubMuted}\itshape #1\par}%
}
\newcommand{\awalan}[1]{%
  {\small\color{ubMuted}#1\par}\vspace{0.18cm}%
}

% --------------------------------------------------------- kotak
\newtcolorbox{kotak}[1][]{enhanced, colback=ubSoft, colframe=ubAccent,
  boxrule=0pt, leftrule=2.2pt, arc=1pt, left=6pt, right=6pt, top=4pt,
  bottom=4pt, fontupper=\footnotesize, #1}
\newtcolorbox{peringatan}[1][]{enhanced, colback=ubWarnBg, colframe=ubWarn,
  boxrule=0pt, leftrule=2.2pt, arc=1pt, left=6pt, right=6pt, top=4pt,
  bottom=4pt, fontupper=\footnotesize, #1}
\newtcolorbox{kiat}[1][]{enhanced, colback=ubGreenBg, colframe=ubGreen,
  boxrule=0pt, leftrule=2.2pt, arc=1pt, left=6pt, right=6pt, top=4pt,
  bottom=4pt, fontupper=\footnotesize, #1}

% --------------------------------------------------------- sampul
% \deckcover{nomor}{minggu}{judul}{sub-cpmk}{teknologi}{sumber}{blok}
\newcommand{\deckcover}[7]{%
\begin{frame}[plain]
\begin{tikzpicture}[remember picture, overlay]
  \fill[ubNavy] (current page.south west) rectangle ([xshift=6.00cm]current page.north west);
  \fill[ubAccent] ([xshift=6.00cm]current page.south west) rectangle ([xshift=6.08cm]current page.north west);

  \node[anchor=north west, text width=4.72cm, align=left, inner sep=0pt]
    at ($(current page.north west)+(0.66cm,-0.60cm)$) {%
      {\color{white}\bfseries\fontsize{10.2}{12.2}\selectfont
       Universitas Bale Bandung}\\[4pt]
      {\color{ubHead}\fontsize{7.6}{9.4}\selectfont
       Fakultas Teknologi Informasi}\\[1pt]
      {\color{ubHead}\fontsize{7.6}{9.4}\selectfont
       Program Studi Teknik Informatika}\\[7pt]
      {\color{ubAccent}\rule{4.0cm}{0.5pt}}%
    };

  \node[anchor=north west, text width=4.72cm, align=left, inner sep=0pt]
    at ($(current page.north west)+(0.66cm,-2.72cm)$) {%
      {\color{ubAccent}\fontsize{6.4}{7.8}\selectfont MATA KULIAH}\\[3pt]
      {\color{white}\bfseries\fontsize{9.8}{12.0}\selectfont
       Pengantar Pemrograman Game}\\[3pt]
      {\color{ubHead}\fontsize{7.6}{9.4}\selectfont
       MWCA009 \textperiodcentered\ 3 SKS}\\[8pt]
      {\color{ubHead}\fontsize{7.6}{9.4}\selectfont
       Semester Ganjil 2026/2027}%
    };

  \node[anchor=south west, text width=4.72cm, align=left, inner sep=0pt]
    at ($(current page.south west)+(0.66cm,0.52cm)$) {%
      {\color{ubAccent}\fontsize{6.4}{7.8}\selectfont DOSEN PENGAMPU}\\[3pt]
      {\color{white}\fontsize{8.4}{10.4}\selectfont
       Cecep Suwanda, S.Si., M.Kom.}%
    };

  \node[anchor=north west, text width=8.95cm, align=left, inner sep=0pt]
    at ($(current page.north east)+(-9.40cm,-0.92cm)$) {%
      {\color{ubAccent}\bfseries\fontsize{8.2}{10.2}\selectfont
       MATERI #1\ \ \textperiodcentered\ \ MINGGU #2}\\[8pt]
      {\color{ubNavy}\bfseries\fontsize{19.4}{23.4}\selectfont
       \hyphenpenalty=10000 \exhyphenpenalty=10000 \raggedright #3}%
    };

  \node[anchor=north west, text width=8.95cm, align=left, inner sep=0pt]
    at ($(current page.north east)+(-9.40cm,-4.30cm)$) {%
      \renewcommand{\arraystretch}{1.30}%
      \begin{tabular}{@{}>{\color{ubMuted}\fontsize{7.3}{9.0}\selectfont}p{1.85cm}
                        >{\color{ubText}\fontsize{8.1}{10.0}\selectfont}p{6.7cm}@{}}
        Sub-CPMK      & #4 \\
        Teknologi     & #5 \\
        Sumber materi & #6 \\
        Blok          & #7 \\
      \end{tabular}%
    };
\end{tikzpicture}
\end{frame}%
}

% --------------------------------------------------------- pembatas bagian
\newcommand{\sectionframe}[3]{%
\begin{frame}[plain]
\begin{tikzpicture}[remember picture, overlay]
  \fill[ubNavy] (current page.south west) rectangle (current page.north east);
  \fill[ubAccent] (current page.north west) rectangle ([xshift=0.10cm]current page.south west);
  \node[anchor=west, inner sep=0pt]
    at ($(current page.west)+(1.00cm,1.45cm)$)
    {\color{ubNavy!52!white}\bfseries\fontsize{58}{58}\selectfont #1};
  \node[anchor=west, inner sep=0pt, text width=11.6cm, align=left]
    at ($(current page.west)+(1.05cm,-0.55cm)$) {%
      {\color{ubAccent}\bfseries\fontsize{8.0}{10.0}\selectfont BAGIAN}\\[3pt]
      {\color{white}\bfseries\fontsize{20.5}{24.5}\selectfont #2}\\[5pt]
      {\color{ubHead}\small #3}%
    };
  \node[anchor=south east, inner sep=0pt]
    at ($(current page.south east)+(-0.60cm,0.32cm)$)
    {\color{ubHead}\fontsize{7.4}{9}\selectfont #1 / 6
     \ \ \textperiodcentered\ \ \insertframenumber};
\end{tikzpicture}
\end{frame}%
}

% --------------------------------------------------------- peta pertemuan
\newcommand{\blokpeta}[6]{%
  \ifnum#2=#3
    \node[rounded corners=3pt, fill=ubNavy, draw=ubNavy, line width=0.5pt,
          text width=3.02cm, minimum height=1.62cm, align=left, inner sep=5pt]
      at (#1,0) {%
        {\color{ubAccent}\bfseries\fontsize{6.8}{8}\selectfont #4}\\[1pt]
        {\color{white}\bfseries\fontsize{7.6}{9}\selectfont #5}\\[2pt]
        {\color{ubHead}\fontsize{6.4}{7.6}\selectfont #6}};
  \else
    \node[rounded corners=3pt, fill=ubSoft, draw=ubHead, line width=0.5pt,
          text width=3.02cm, minimum height=1.62cm, align=left, inner sep=5pt]
      at (#1,0) {%
        {\color{ubAccent}\bfseries\fontsize{6.8}{8}\selectfont #4}\\[1pt]
        {\color{ubNavy}\bfseries\fontsize{7.6}{9}\selectfont #5}\\[2pt]
        {\color{ubMuted}\fontsize{6.4}{7.6}\selectfont #6}};
  \fi
}

% \petapertemuan{blok aktif}{sebelumnya}{sesudahnya}{keluaran}
% \petapertemuan{blok aktif}{sebelumnya}{pertemuan ini}{sesudahnya}{keluaran}
\newcommand{\petapertemuan}[5]{%
\begin{frame}{Peta Pertemuan}
{\scriptsize\color{ubMuted}Kedudukan materi dalam empat blok perkuliahan}\par
\vspace{0.10cm}
\centering
\begin{tikzpicture}
  \blokpeta{-5.58}{1}{#1}{BLOK I}{Materi 1--5}{Game loop, vektor,\\fisika, collision}
  \blokpeta{-1.86}{2}{#1}{BLOK II}{Materi 6--9}{AI NPC: chase, FSM,\\pathfinding A*}
  \blokpeta{1.86}{3}{#1}{BLOK III}{Materi 10--11}{Unity: scene,\\level, physics}
  \blokpeta{5.58}{4}{#1}{BLOK IV}{Materi 12--14}{Aset visual, audio,\\UI, playtest}
\end{tikzpicture}
\par\vspace{0.30cm}
\begin{columns}[T,onlytextwidth]
  \begin{column}{0.315\textwidth}
    \begin{beamercolorbox}[wd=\linewidth,sep=4pt,rounded=true]{block body}
      {\color{ubMuted}\fontsize{6.6}{8}\selectfont SEBELUMNYA}\\[2pt]
      {\footnotesize #2}
    \end{beamercolorbox}
  \end{column}
  \begin{column}{0.315\textwidth}
    \begin{beamercolorbox}[wd=\linewidth,sep=4pt,rounded=true]{block title}
      {\color{ubHead}\fontsize{6.6}{8}\selectfont PERTEMUAN INI}\\[2pt]
      {\footnotesize #3}
    \end{beamercolorbox}
  \end{column}
  \begin{column}{0.315\textwidth}
    \begin{beamercolorbox}[wd=\linewidth,sep=4pt,rounded=true]{block body}
      {\color{ubMuted}\fontsize{6.6}{8}\selectfont SESUDAHNYA}\\[2pt]
      {\footnotesize #4}
    \end{beamercolorbox}
  \end{column}
\end{columns}
\par\vspace{0.26cm}
{\footnotesize\color{ubText}\textbf{Keluaran pertemuan.} #5\par}
\end{frame}%
}

% --------------------------------------------------------- penutup
\newcommand{\penutup}[2]{%
\begin{frame}[plain]
\begin{tikzpicture}[remember picture, overlay]
  \fill[ubNavy] (current page.south west) rectangle (current page.north east);
  \fill[ubAccent] (current page.north west) rectangle ([xshift=0.10cm]current page.south west);

  \node[anchor=north west, text width=13.6cm, align=left, inner sep=0pt]
    at ($(current page.north west)+(1.00cm,-0.95cm)$) {%
      {\color{ubAccent}\bfseries\fontsize{8.2}{10}\selectfont ASESMEN MINGGU INI}\\[6pt]
      {\color{white}\bfseries\fontsize{24}{28}\selectfont #1}%
    };

  \node[anchor=north west, text width=13.6cm, align=left, inner sep=0pt]
    at ($(current page.north west)+(1.00cm,-4.55cm)$) {%
      {\color{ubNavy!0!ubHead}\rule{13.4cm}{0.4pt}}\\[7pt]
      {\color{ubAccent}\bfseries\fontsize{7.4}{9}\selectfont
       PENGAIT KE MATERI BERIKUTNYA}\\[4pt]
      {\color{ubHead}\fontsize{10.6}{13.2}\selectfont #2}%
    };

  \node[anchor=south east, inner sep=0pt]
    at ($(current page.south east)+(-0.60cm,0.32cm)$)
    {\color{ubHead}\fontsize{7.4}{9}\selectfont \insertframenumber};
\end{tikzpicture}
\end{frame}%
}



\renewcommand{\materino}{14}

\begin{document}

\deckcover{14}{15}
  {Evaluasi Prototipe, Playtesting, dan Presentasi}
  {4.3}
  {Unity}
  {Bab XVI buku ajar}
  {Blok IV \textperiodcentered\ Materi 12--14}

% --------------------------------------------------------- capaian
\begin{frame}{Capaian Pertemuan}
\setbeamerfont{itemize/enumerate body}{size=\scriptsize}
\awalan{Sub-CPMK dan indikator dari RPS, serta tujuan Bab XVI buku ajar. Bab ini
menutup \emph{prototipe} sekaligus menutup seluruh rangkaian materi.}
\begin{columns}[T,onlytextwidth]
  \begin{column}{0.52\textwidth}
    \begin{beamercolorbox}[wd=\linewidth,sep=5pt,rounded=true]{block title}
      {\fontsize{6.8}{8}\selectfont\color{ubHead}SUB-CPMK 4.3 \textperiodcentered\ RANAH C5}\\[2pt]
      {\footnotesize Mengevaluasi prototipe \emph{game} melalui
       \emph{playtesting} dan memperbaiki integrasi level serta aset.}
    \end{beamercolorbox}
    \vspace{0.16cm}

    {\fontsize{7}{8.5}\selectfont\bfseries\color{ubAccent}INDIKATOR}\par
    \vspace{0.06cm}
    \begin{itemize}
      \item Menyusun \emph{checklist fungsional} dan lembar \emph{playtest}.
      \item Mengidentifikasi \emph{bug} dan masalah integrasi dari hasil
            \emph{playtest}.
      \item Menyajikan laporan perbaikan dan refleksi pemilihan teknologi.
    \end{itemize}
  \end{column}
  \begin{column}{0.45\textwidth}
    {\fontsize{7}{8.5}\selectfont\bfseries\color{ubAccent}TUJUAN BAB XVI}\par
    \vspace{0.08cm}
    \begin{itemize}
      \item Menjelaskan mengapa pembuat tidak dapat menilai karyanya sendiri.
      \item Menyusun \emph{checklist fungsional} dari GDD sendiri.
      \item Menjalankan pemeriksaan otomatis untuk \emph{soft-lock} dan jalan
            yang hanya menurun.
      \item Menyusun lembar \emph{playtest} dan melaksanakannya antar
            kelompok.
      \item Mencatat \emph{bug} yang dapat diulang dan menimbang beratnya.
      \item Menyajikan demo 5--10 menit dan menyusun refleksi teknologi.
    \end{itemize}
  \end{column}
\end{columns}
\end{frame}

% --------------------------------------------------------- peta
\petapertemuan{4}
  {Materi 13 \textperiodcentered\ Integrasi Audio dan \emph{Trigger} Interaksi}
  {Materi 14 \textperiodcentered\ Evaluasi Prototipe, \emph{Playtesting}, dan
   Presentasi}
  {UAS (Minggu 16)}
  {Prototipe teruji: \emph{checklist} fungsional terisi, \emph{log bug} beserta
   status penanganannya, demo 5--10 menit, laporan perbaikan, dan refleksi
   teknologi.}

% ========================================================= bagian 1
\section{Mengapa Prototipe Harus Diuji}
\sectionframe{01}{Mengapa Prototipe Harus Diuji}
  {Pembuat tidak dapat menilai karyanya sendiri}

\begin{frame}{Pembuat Tidak Dapat Menilai Karyanya Sendiri}
\awalan{Selama ini Anda menguji permainan dengan satu cara: menjalankannya,
lalu memeriksa apakah perilakunya sesuai maksud Anda. Cara itu gagal pada hal
yang paling penting.}
\begin{itemize}
  \item Anda sudah tahu tombol mana yang harus ditekan, jalur mana yang harus
        dilewati, dan ke mana tujuan yang sebenarnya.
  \item Pengetahuan itu tidak dimiliki pemain baru, dan justru pengetahuan
        itulah yang menutupi masalah.
  \item Level yang tampak jelas bagi pembuatnya dapat membingungkan pemain
        lain selama beberapa menit pertama.
  \item Karena itu pengujian menuntut orang lain yang belum pernah melihat
        karya Anda, dan menuntut Anda menahan diri untuk tidak menolongnya.
\end{itemize}
\vspace{0.12cm}
\begin{peringatan}
Kesalahan yang paling sering terjadi pada sesi \emph{playtest} pertama adalah
memberi petunjuk ketika penguji kebingungan. Kalimat ``tekan tombol spasi untuk
lompat'' tampak seperti bantuan, padahal ia menghapus satu-satunya keterangan
yang paling Anda butuhkan. Apabila penguji tidak menemukan sendiri caranya, itu
temuan, bukan kesalahan penguji.
\end{peringatan}
\end{frame}

\begin{frame}{Dua Jenis Pengujian yang Saling Melengkapi}
\setbeamerfont{itemize/enumerate body}{size=\scriptsize}
\awalan{Keduanya menemukan hal yang berbeda, dan urutannya penting.}
\begin{columns}[T,onlytextwidth]
  \begin{column}{0.48\textwidth}
    \begin{block}{Pemeriksaan otomatis}
      \begin{itemize}
        \item Dijalankan komputer, tanpa pemain, dan berulang-ulang.
        \item Menemukan hal yang bersifat pasti: petak yang tidak dapat
              dicapai, lompatan yang melebihi jangkauan, berkas yang gagal
              dimuat.
        \item Cepat, dan hasilnya sama setiap kali dijalankan.
      \end{itemize}
    \end{block}
  \end{column}
  \begin{column}{0.48\textwidth}
    \begin{exampleblock}{\emph{Playtest}}
      \begin{itemize}
        \item Dijalankan oleh orang lain yang memainkan prototipe Anda.
        \item Menemukan hal yang tidak dapat dihitung: bagian yang
              membingungkan, membosankan, dan membuat pemain berhenti.
        \item Menuntut penguji yang belum pernah melihat karya Anda.
      \end{itemize}
    \end{exampleblock}
  \end{column}
\end{columns}
\vspace{0.16cm}
\begin{kotak}
Pemeriksaan otomatis dijalankan lebih dahulu, sebab tidak ada gunanya meminta
orang lain menguji level yang belum dapat diselesaikan. Namun ia tidak dapat
menggantikan \emph{playtest}: mesin dapat memberitahu bahwa levelnya
\emph{mungkin} diselesaikan, tetapi tidak dapat memberitahu apakah levelnya
\emph{menyenangkan} diselesaikan.
\end{kotak}
\end{frame}

\begin{frame}{Pengujian sebagai Siklus, Bukan Langkah Terakhir}
\centering
\adjustbox{max width=0.98\linewidth,max height=0.55\textheight}{%
  % chapters/16/figures/siklus-uji.tikz
% Pengujian sebagai siklus: setiap perbaikan menuntut pengujian ulang.
\begin{tikzpicture}[
  tahap/.style={
    rectangle, rounded corners=5pt, draw=ubPrimary, thick, fill=ubAlt,
    align=center, text width=27mm, inner sep=4pt, minimum height=11mm,
    font=\footnotesize
  },
  alat/.style={
    rectangle, rounded corners=5pt, draw=ubAccent, thick, fill=white,
    align=center, text width=27mm, inner sep=4pt, minimum height=11mm,
    font=\footnotesize
  },
  panah/.style={-{Stealth[length=2.2mm]}, thick, draw=ubAccent},
  kembali/.style={-{Stealth[length=2.2mm]}, thick, draw=ubPrimary, dashed},
  catatan/.style={font=\scriptsize\itshape, align=left}
]

\node[tahap] (gdd)    at (0,0)     {\istilah{GDD}\\{\scriptsize\itshape rencana}};
\node[alat]  (daftar) at (3.9,0)   {checklist\\fungsional};
\node[alat]  (otomat) at (7.8,0)   {uji otomatis\\{\scriptsize\itshape tanpa pemain}};
\node[alat]  (main)   at (11.7,0)  {playtest\\{\scriptsize\itshape dengan pemain}};

\node[tahap] (log)    at (11.7,-2.0) {log \istilah{bug}};
\node[tahap] (perbaiki) at (7.8,-2.0) {perbaikan};
\node[tahap] (build)  at (3.9,-2.0) {build akhir};

\draw[panah] (gdd) -- (daftar);
\draw[panah] (daftar) -- (otomat);
\draw[panah] (otomat) -- (main);
\draw[panah] (main) -- (log);
\draw[panah] (log) -- (perbaiki);
\draw[panah] (perbaiki) -- (build);

% Kembali ke awal: perbaikan selalu menuntut pengujian ulang.
\draw[kembali] (perbaiki.north) .. controls (7.8,1.9) and (5.6,1.9) .. (daftar.north);

\node[catatan, anchor=north west, text width=62mm] at (-1.6,-2.9)
  {Perbaikan tidak mengakhiri pekerjaan, melainkan memulai putaran
   berikutnya. Garis putus-putus menandai bahwa setiap perbaikan menuntut
   pemeriksaan ulang, dan pemeriksaan ulang menuntut checklist yang
   diperbarui.};
\end{tikzpicture}
}\par
\vspace{0.12cm}
\begin{kotak}
Perbaikan atas satu temuan berpotensi menimbulkan temuan baru: menambah pijakan
agar pemain tidak terjebak dapat membuat bagian lain terlalu mudah, dan
menggeser letak tujuan dapat membuat satu bagian tidak lagi terjangkau. Karena
itu siklusnya kembali ke daftar periksa, bukan berhenti pada \emph{build} akhir.
\end{kotak}
\end{frame}

% ========================================================= bagian 2
\section{Checklist Fungsional}
\sectionframe{02}{Checklist Fungsional}
  {Daftar janji yang diturunkan dari GDD}

\begin{frame}{Menurunkan Checklist dari GDD}
\awalan{\emph{Checklist} fungsional bukan daftar keinginan, melainkan daftar
janji: setiap butirnya berasal dari sesuatu yang sudah Anda tuliskan.}
\begin{itemize}
  \item \emph{Checklist} fungsional adalah daftar pernyataan yang harus
        bernilai benar sebelum \emph{build} dianggap selesai.
  \item Sumbernya adalah \emph{Game Design Document} (GDD), yaitu rancangan
        tertulis yang Anda susun pada awal pengerjaan proyek.
  \item Cara menyusunnya tidak menuntut pekerjaan baru, hanya ketelitian
        membaca kembali: buka GDD, lalu ubah setiap pernyataan menjadi satu
        butir yang dapat diperiksa benar atau salah.
  \item Pernyataan GDD bernilai rancangan, sedangkan butir \emph{checklist}
        bernilai pemeriksaan. Di situlah gunanya: niat berubah menjadi syarat
        yang dapat diuji, termasuk oleh program.
\end{itemize}
\vspace{0.14cm}
\begin{kiat}
Butir yang baik dapat dijawab ``ya'' atau ``tidak'' oleh orang lain tanpa
bertanya kepada Anda. Butir yang masih memuat kata ``menarik'' atau ``cukup''
belum dapat diperiksa, dan karena itu belum berguna.
\end{kiat}
\end{frame}

\begin{frame}{Tiga Kelompok Butir}
\awalan{\emph{Checklist} yang hanya memuat satu jenis butir akan meninggalkan
bagian yang penting.}
\begin{enumerate}
  \item \textbf{Mekanik.} Apakah setiap mekanik yang dijanjikan GDD benar-benar
        ada dan berjalan? Lompatan, pengambilan benda, kondisi menang, dan
        kondisi kalah.
  \item \textbf{Integrasi.} Apakah bagian-bagian yang dibuat pada bab yang
        berbeda benar-benar tersambung? Bunyi terdengar pada peristiwanya,
        animasi berganti menurut keadaannya, antarmuka menampilkan angka yang
        benar.
  \item \textbf{Keseluruhan.} Apakah permainannya dapat diselesaikan dari awal
        sampai akhir, tanpa berhenti, tanpa jalan buntu, dan dalam waktu yang
        Anda janjikan?
\end{enumerate}
\vspace{0.14cm}
\begin{kotak}
Kelompok ketiga paling sering terlewat, sebab ia tidak dapat diperiksa sambil
membuat. Ia hanya dapat diperiksa dengan menjalankan permainan dari awal sampai
akhir, dan itulah sebabnya ia memerlukan \emph{build}, bukan \emph{scene} yang
sedang Anda sunting.
\end{kotak}
\end{frame}

\begin{frame}{Checklist sebagai Berkas}
\awalan{Simpan \emph{checklist} sebagai berkas, bukan hafalan atau ingatan,
sehingga hasil setiap putaran dapat dibandingkan.}
{\scriptsize
\renewcommand{\arraystretch}{1.12}
\begin{tabularx}{\textwidth}{@{}>{\raggedright\arraybackslash}p{0.55cm}
    >{\raggedright\arraybackslash}p{1.50cm}
    >{\raggedright\arraybackslash}X
    >{\raggedright\arraybackslash}p{2.00cm}@{}}
\toprule
\textbf{No} & \textbf{Kelompok} & \textbf{Butir pemeriksaan} &
  \textbf{Putaran 1/2} \\
\midrule
1 & Mekanik & Pemain dapat melompat hanya dari tanah; tidak ada lompatan di
  udara. & Lulus / Lulus \\
2 & Mekanik & Tiga benda \texttt{Pickup} tersedia dan pintu akhir terbuka
  setelah ketiganya diambil. & Gagal / Lulus \\
3 & Integrasi & Bunyi pengambilan terdengar pada setiap benda, termasuk benda
  ketiga. & Gagal / Lulus \\
4 & Integrasi & Animasi diam dan berlari berganti sesuai gerak pemain. &
  Lulus / Lulus \\
5 & Keseluruhan & Satu penyelesaian penuh berlangsung 5--10 menit tanpa jalan
  buntu. & Gagal / Lulus \\
\bottomrule
\end{tabularx}}
\vspace{0.12cm}
\begin{itemize}
  \item Batasi pada sekitar dua puluh sampai tiga puluh butir: daftar yang
        terlalu panjang tidak akan pernah dijalankan sampai selesai, sedangkan
        yang terlalu pendek melewatkan hal penting.
  \item Butir yang paling berguna adalah butir yang pernah gagal pada putaran
        sebelumnya.
\end{itemize}
\end{frame}

% ========================================================= bagian 3
\section{Pemeriksaan Otomatis Tanpa Pemain}
\sectionframe{03}{Pemeriksaan Otomatis Tanpa Pemain}
  {\emph{Soft-lock}, jangkauan lompatan, dan pencarian jalur}

\begin{frame}{\emph{Soft-lock}: Kegagalan yang Tidak Terlihat}
\awalan{Satu jenis kesalahan tidak menimbulkan pesan apa pun, tidak
menghentikan program, dan tidak terlihat pada satu kali percobaan.}
\begin{itemize}
  \item \emph{Soft-lock} adalah keadaan ketika pemain tidak dapat melanjutkan
        dan tidak dapat kembali, walaupun permainannya sendiri masih berjalan.
  \item Program tidak berhenti, jendela \emph{Console} tidak mencetak kesalahan,
        dan penguji yang mencoba satu kali mungkin tidak menemukannya.
  \item Contohnya pemain terjatuh ke dalam lubang yang terlalu tinggi untuk
        ditinggalkan.
  \item Hanya sebagian pemain mengalaminya, dan itulah yang membuat kesalahan
        ini mahal.
\end{itemize}
\vspace{0.14cm}
\begin{kotak}
Karena kesalahan ini bersifat pasti, bukan soal selera, ia dapat diperiksa
secara pasti pula: periksa setiap tempat pemain dapat berdiri, lalu tanyakan
apakah dari sana titik tujuan masih dapat dicapai.
\end{kotak}
\end{frame}

\begin{frame}[fragile]{Jangkauan Lompatan, Dihitung lalu Diputuskan}
\lstinputlisting[language=Python,firstline=65,lastline=78]
  {../buku_ajar/code/python/ch16_uji_mandiri.py}
\kodejelas{Dengan $v_0=700$ piksel per detik, $g=1800$ piksel per detik
kuadrat, dan kecepatan jalan 300 piksel per detik, tinggi lompatan 136,1 piksel
(4,25 petak) dan jarak mendatar 233,3 piksel (7,29 petak) pada petak 32 piksel.
Melompat ke tempat yang lebih rendah memberi waktu di udara lebih lama,
sehingga jangkauannya lebih jauh; karena itu jangkauan dihitung untuk setiap
pasangan tempat berdiri.}
\end{frame}

\begin{frame}[fragile]{Memeriksa Setiap Tempat Berdiri}
\lstinputlisting[language=Python,firstline=277,lastline=286]
  {../buku_ajar/code/python/ch16_uji_mandiri.py}
\kodejelas{Setiap pijakan memanggil \texttt{cari\_jalur}; pijakan yang tidak
menemukan jalur dicatat sebagai calon \emph{soft-lock}. Algoritmanya sama
dengan Bab X, dan yang berbeda hanya aturan tetangganya: berjalan, jatuh, atau
melompat.}
\vspace{0.12cm}
\begin{kotak}
Pada peta ketiga, simulasi 200 titik awal berhasil 82,0\%,
sedangkan pemeriksaan menyeluruh menemukan 15,6\% tempat terjebak. Simulasi
hanya menyentuh keadaan yang kebetulan terpilih; yang menyingkap
\emph{soft-lock} adalah pemeriksaan atas \emph{setiap} tempat berdiri.
\end{kotak}
\end{frame}

% ========================================================= bagian 4
\section{Lembar dan Pelaksanaan Playtest}
\sectionframe{04}{Lembar dan Pelaksanaan \emph{Playtest}}
  {Mengamati, mencatat, dan bertukar peran}

\begin{frame}{Mengamati, Bukan Bertanya}
\awalan{Pertanyaan ``apakah permainannya bagus'' tidak berguna, sebab
jawabannya sopan dan tidak dapat ditindaklanjuti.}
\begin{itemize}
  \item Yang berguna adalah pengamatan atas apa yang benar-benar dilakukan
        penguji, bukan pendapatnya tentang karya Anda.
  \item Ketika penguji berhenti dua puluh detik di depan pintu yang belum
        terbuka lalu berjalan ke arah yang salah, Anda memperoleh satu temuan
        yang dapat langsung diperbaiki.
  \item Catat waktu, bukan hanya kejadian: ``berhenti 20 detik'' jauh lebih
        berguna daripada ``sempat bingung''.
  \item Angka waktu memudahkan Anda membandingkan hasil perbaikan dengan hasil
        sebelumnya.
\end{itemize}
\vspace{0.14cm}
\begin{kiat}
Menit pertama adalah menit yang paling banyak memuat temuan, dan menit itu pula
yang paling cepat Anda lupakan. Amati dan catat sejak awal, lalu kumpulkan
kesan sesudah permainan selesai, bukan di tengah jalan.
\end{kiat}
\end{frame}

\begin{frame}{Isi Lembar \emph{Playtest}}
\awalan{Lembar \emph{playtest} memuat keterangan tentang penguji dan catatan
atas apa yang diamati selama permainan berlangsung.}
{\scriptsize
\renewcommand{\arraystretch}{1.12}
\begin{tabularx}{\textwidth}{@{}>{\raggedright\arraybackslash}p{3.5cm}
    >{\raggedright\arraybackslash}X@{}}
\toprule
\textbf{Bagian} & \textbf{Isi} \\
\midrule
Keterangan penguji & Nama atau kelompok, dan pengalaman bermain permainan
  sejenis. \\
\addlinespace
Berkas yang diuji & Nama \emph{build} beserta tanggalnya. \\
\addlinespace
Waktu penyelesaian & Lama permainan dari awal sampai tujuan, atau keterangan
  bahwa pemain berhenti sebelum selesai. \\
\addlinespace
Titik berhenti & Tempat dan waktu ketika penguji berhenti lebih lama dari
  sepuluh detik, beserta apa yang tampak ia cari. \\
\addlinespace
Kesalahan yang terlihat & Setiap keadaan yang tidak seharusnya terjadi,
  beserta urutan tindakan yang menyebabkannya. \\
\addlinespace
Masukan penguji & Pendapat penguji, dicatat apa adanya, tanpa dibantah. \\
\addlinespace
Catatan penguji sendiri & Hal yang Anda amati tetapi tidak disebutkan
  penguji. \\
\bottomrule
\end{tabularx}}
\end{frame}

\begin{frame}{\emph{Playtest} antar Kelompok: Enam Langkah}
\awalan{Keuntungannya berjalan dua arah: kelompok Anda memperoleh penguji yang
belum pernah melihat karya Anda, dan Anda memperoleh pengalaman menguji karya
orang lain.}
\begin{enumerate}
  \item \textbf{Siapkan \emph{build} yang dapat dijalankan} pada komputer lain,
        bukan pada komputer Anda. Berkas yang hanya berjalan di komputer
        pembuatnya bukan \emph{build}.
  \item \textbf{Tuliskan tujuan sesi}: bagian mana yang paling Anda ragukan.
        Pengujian yang diarahkan menemukan lebih banyak daripada yang
        dibiarkan bebas.
  \item \textbf{Serahkan kendali tanpa memberi petunjuk.} Jawablah pertanyaan
        dengan ``apa yang Anda duga seharusnya terjadi''.
  \item \textbf{Amati dan catat}, terutama menit pertama; kumpulkan kesan
        sesudah permainan selesai.
  \item \textbf{Tukar peran.} Uji karya kelompok lain dengan cara yang sama,
        agar Anda memahami bagaimana rasanya diamati.
  \item \textbf{Ringkas temuan segera}, lalu pindahkan temuan yang dapat diulang
        ke \emph{log bug}.
\end{enumerate}
\begin{peringatan}
Jangan membela karya Anda selama sesi berlangsung. Penguji yang merasa
pendapatnya diperdebatkan akan berhenti memberi masukan, dan sesi berikutnya
menjadi sia-sia.
\end{peringatan}
\end{frame}

% ========================================================= bagian 5
\section{Log Bug dan Perbaikan}
\sectionframe{05}{Log \emph{Bug} dan Perbaikan}
  {Catatan yang dapat diulang dan beratnya}

\begin{frame}{Catatan yang Dapat Diulang}
\awalan{``Bunyinya kadang tidak terdengar'' tidak dapat diperbaiki, sebab Anda
tidak tahu kapan ia terjadi. Setiap catatan harus memuat tiga hal.}
\begin{itemize}
  \item Apa yang terjadi, apa yang seharusnya terjadi, dan urutan tindakan yang
        membuatnya terjadi.
  \item Catatan baru boleh dipindahkan ke \emph{log bug} setelah ada urutan
        yang dapat diulang oleh orang lain.
\end{itemize}
\vspace{0.10cm}
{\scriptsize
\renewcommand{\arraystretch}{1.12}
\begin{tabularx}{\textwidth}{@{}>{\raggedright\arraybackslash}p{0.55cm}
    >{\raggedright\arraybackslash}p{1.55cm}
    >{\raggedright\arraybackslash}X
    >{\raggedright\arraybackslash}p{1.15cm}
    >{\raggedright\arraybackslash}p{1.10cm}@{}}
\toprule
\textbf{No} & \textbf{Jenis} & \textbf{Urutan penyebab dan hasil yang salah} &
  \textbf{Berat} & \textbf{Status} \\
\midrule
1 & \emph{Soft-lock} & Pemain berjalan ke kiri di ujung lorong, lalu jatuh ke
  lubang. Seluruh tombol tidak lagi mengeluarkan pemain; permainan harus dimulai
  ulang. & Berat & Belum \\
\addlinespace
2 & Integrasi & Bunyi pengambilan tidak terdengar pada benda ketiga. Benda
  pertama dan kedua berbunyi; skornya tetap bertambah. & Sedang & Belum \\
\addlinespace
3 & Tampilan & Animasi diam tetap berjalan ketika pemain berlari, apabila
  berlari dimulai ketika pemain baru saja mendarat. & Ringan & Belum \\
\bottomrule
\end{tabularx}}
\end{frame}

\begin{frame}{Menimbang Beratnya}
\awalan{Tidak semua temuan perlu diperbaiki, dan waktu menjelang pengumpulan
selalu terbatas. Setiap temuan ditimbang atas dua pertanyaan.}
\begin{enumerate}
  \item \textbf{Apakah permainannya masih dapat diselesaikan?} Temuan yang
        menjawab ``tidak'' selalu paling berat, walaupun ia jarang terjadi.
  \item \textbf{Berapa banyak pemain yang akan menemukannya?} Temuan yang
        dialami setiap pemain pada menit pertama mendesak; temuan yang hanya
        muncul pada keadaan jarang dapat ditunda.
\end{enumerate}
\vspace{0.10cm}
{\scriptsize
\renewcommand{\arraystretch}{1.12}
\begin{tabularx}{\textwidth}{@{}>{\raggedright\arraybackslash}p{1.55cm}
    >{\raggedright\arraybackslash}p{5.20cm}
    >{\raggedright\arraybackslash}X@{}}
\toprule
\textbf{Berat} & \textbf{Tandanya} & \textbf{Tindakan} \\
\midrule
Berat & Permainan tidak dapat diselesaikan, atau kemajuannya hilang. &
  Diperbaiki sebelum sesi \emph{playtest} berikutnya. \\
\addlinespace
Sedang & Permainan selesai, tetapi satu bagian yang dijanjikan GDD tidak
  bekerja. & Diperbaiki, dan diperiksa kembali pada putaran berikutnya. \\
\addlinespace
Ringan & Tampilan atau bunyi yang tidak sesuai, tanpa mengubah jalannya
  permainan. & Diperbaiki apabila waktu tersisa; dicatat sebagai kekurangan
  yang diketahui. \\
\bottomrule
\end{tabularx}}
\begin{kiat}
Perbaiki temuan berat satu per satu, lalu jalankan \emph{checklist} fungsional
sesudah setiap perbaikan. Menggabungkan beberapa perbaikan membuat Anda tidak
tahu mana yang bekerja.
\end{kiat}
\end{frame}

\begin{frame}{Dari Temuan kembali ke GDD}
\awalan{Lingkaran pengujian belum selesai pada perbaikan. Temuan yang berulang
perlu diperiksa terhadap rancangannya.}
\begin{itemize}
  \item Temuan yang berulang pada beberapa sesi biasanya bukan kesalahan
        pelaksanaan, melainkan tanda bahwa rancangan pada GDD memang belum
        jelas.
  \item Apabila setiap penguji melewati pintu yang belum terbuka, kemungkinan
        besar GDD tidak pernah menjelaskan bahwa pintu itu terkunci, sehingga
        petunjuknya juga tidak pernah dipasang.
  \item Temuan yang berulang diperiksa terhadap GDD, lalu GDD diperbarui.
  \item Dengan begitu siklusnya benar-benar kembali ke awal: perbaikan menuntut
        \emph{checklist} yang diperbarui, dan \emph{checklist} yang diperbarui
        menuntut pemeriksaan ulang.
\end{itemize}
\vspace{0.14cm}
\begin{kotak}
Jangan menghapus janji dari GDD agar hasilnya tampak sesuai. Itu bentuk
ketidakjujuran akademik, dan sekaligus menghapus bukti bahwa Anda mampu menilai
karya sendiri.
\end{kotak}
\end{frame}

% ========================================================= bagian 6
\section{Presentasi Demo, Kesesuaian GDD, dan Refleksi Teknologi}
\sectionframe{06}{Presentasi Demo dan Refleksi Teknologi}
  {Demo, kesesuaian rancangan, dan pilihan perkakas}

\begin{frame}{Demo Lima sampai Sepuluh Menit}
\awalan{Waktu sependek itu menuntut pilihan. Kesalahan yang paling sering
adalah menceritakan seluruh proses pembuatan; yang perlu diperlihatkan bukan
perjalanan Anda, melainkan hasilnya.}
{\scriptsize
\renewcommand{\arraystretch}{1.12}
\begin{tabularx}{\textwidth}{@{}>{\raggedright\arraybackslash}p{1.55cm}
    >{\raggedright\arraybackslash}p{2.60cm}
    >{\raggedright\arraybackslash}X@{}}
\toprule
\textbf{Waktu} & \textbf{Bagian} & \textbf{Isi} \\
\midrule
1 menit & Pembukaan & Nama permainan, jenisnya, dan satu kalimat tentang apa
  yang pemain lakukan di dalamnya. \\
\addlinespace
4--6 menit & Demo langsung & Permainan dijalankan dari awal sampai tujuan, atau
  sampai bagian yang paling Anda banggakan. Tidak ada penjelasan bertele-tele
  di tengah demo. \\
\addlinespace
1--2 menit & Temuan pengujian & Satu atau dua temuan terpenting dari
  \emph{playtest} antar kelompok, beserta perbaikan yang Anda lakukan. \\
\addlinespace
1 menit & Penutup & Kekurangan yang Anda ketahui, dan satu hal yang akan Anda
  kerjakan apabila waktunya bertambah. \\
\bottomrule
\end{tabularx}}
\vspace{0.08cm}
\begin{kiat}
Siapkan demo cadangan: rekam permainan sebagai video pendek, dan bawa salinan
\emph{build} pada dua tempat. Kegagalan teknis saat presentasi adalah hal yang
biasa, dan cara menghadapinya dinilai.
\end{kiat}
\end{frame}

\begin{frame}{Kesesuaian GDD dengan \emph{Build}}
\awalan{GDD ditulis ketika sebagian besar hal belum diketahui. Perbedaan antara
GDD dan \emph{build} hampir pasti ada; yang salah adalah perbedaan yang tidak
Anda sadari.}
{\scriptsize
\renewcommand{\arraystretch}{1.12}
\begin{tabularx}{\textwidth}{@{}>{\raggedright\arraybackslash}p{2.75cm}
    >{\raggedright\arraybackslash}X
    >{\raggedright\arraybackslash}X@{}}
\toprule
\textbf{Keadaan} & \textbf{Artinya} & \textbf{Tindakan} \\
\midrule
Sesuai & Yang dijanjikan GDD benar-benar ada di dalam \emph{build}. &
  Tidak ada yang perlu dikerjakan; cantumkan sebagai bukti pada laporan. \\
\addlinespace
GDD diperbarui & Rancangannya berubah ketika dikerjakan, dan perubahan itu
  disengaja. & Perbarui GDD dan tuliskan alasan perubahannya. Yang dinilai
  adalah kesadaran Anda atas perubahan itu. \\
\addlinespace
Belum dikerjakan & Yang dijanjikan GDD belum ada, dan tidak akan ada. &
  Cantumkan sebagai kekurangan yang diketahui. Jangan menghapusnya dari GDD
  agar laporannya tampak lengkap. \\
\bottomrule
\end{tabularx}}
\vspace{0.12cm}
\begin{itemize}
  \item Baris ketiga adalah yang paling penting untuk dihindari, sebab
        menghapus janji dari GDD menghilangkan bukti kemampuan menilai karya
        sendiri.
\end{itemize}
\end{frame}

\begin{frame}{Refleksi Pemilihan Teknologi}
\awalan{Refleksi bukan penilaian mana yang lebih baik, melainkan keterangan
tentang mana yang cocok untuk keperluan apa.}
\begin{itemize}
  \item \textbf{Pada bagian mana Pygame mempercepat pekerjaan Anda?} Jawaban
        yang umum: bagian yang perhitungannya perlu Anda kuasai sendiri, dan
        permainan yang aturannya sederhana.
  \item \textbf{Pada bagian mana Unity mempercepat pekerjaan Anda?} Jawaban
        yang umum: penyusunan level, pengaturan fisika, animasi, dan bunyi,
        sebab pekerjaan itu diurus mesin.
  \item \textbf{Bagian mana yang paling Anda kuasai sekarang, dan bagian mana
        yang masih perlu dipelajari?} Jawaban atas pertanyaan ini menjadi
        rencana belajar Anda sesudah mata kuliah ini berakhir.
  \item Setiap jawaban sebaiknya disertai contoh dari pekerjaan kelompok Anda
        sendiri, bukan dari buku ini.
\end{itemize}
\vspace{0.14cm}
\begin{kotak}
Refleksi yang jujur akan memuat hal yang tidak Anda kuasai. Menyebutkannya bukan
pengakuan kelemahan, melainkan bukti bahwa Anda mampu menilai kemampuan sendiri
dengan tepat. Itulah bagian tertinggi ranah C5 yang dinilai pada Sub-CPMK 4.3:
menilai berdasarkan bukti, bukan berdasarkan kesan.
\end{kotak}
\end{frame}

% --------------------------------------------------------- ketentuan proyek
\begin{frame}{Ketentuan Proyek Akhir (Minggu 15)}
\setbeamerfont{itemize/enumerate body}{size=\scriptsize}
\awalan{Seluruh ketentuan berikut diambil dari RPS dan diperiksa pada presentasi
Minggu 15.}
\begin{itemize}
  \item \textbf{Persiapan pengujian.} \emph{Checklist} fungsional dan lembar
        \emph{playtest} terstruktur telah disusun sebelum sesi.
  \item \textbf{\emph{Playtest} antar kelompok} dilaksanakan pada Minggu 15;
        seluruh \emph{bug} yang ditemukan dicatat pada \emph{log bug} beserta
        status penanganannya.
  \item \textbf{Kesesuaian GDD dengan \emph{build}} yang benar-benar dijalankan
        diperiksa dan dilaporkan, termasuk perbedaan yang disengaja.
  \item \textbf{Demo 5--10 menit} disajikan menurut susunan pembukaan, demo
        langsung, temuan pengujian, dan penutup.
  \item \textbf{Refleksi pemilihan teknologi}: Python/Pygame dibandingkan
        dengan Unity, disertai contoh dari pekerjaan sendiri.
  \item \textbf{Keluaran proyek}: \emph{build}, \emph{project source}, ringkasan
        desain, laporan pengembangan, dan lembar \emph{playtest} terisi; uji
        fungsional oleh dosen pada Minggu 16 (UAS).
\end{itemize}
\vspace{0.12cm}
\begin{kotak}
Yang dinilai bukan permainan yang tampak tanpa cacat, melainkan kemampuan
menemukan masalah pada karya sendiri, memperbaikinya, dan
mempertanggungjawabkan pilihan teknologi.
\end{kotak}
\end{frame}

% --------------------------------------------------------- aktivitas
\begin{frame}{Aktivitas di Kelas \textperiodcentered\ Playtest}
\setbeamerfont{itemize/enumerate body}{size=\scriptsize}
\awalan{Pengalaman belajar minggu 15 pada RPS, dikerjakan berurutan di kelas
sebagai bagian pengujian Proyek Akhir.}
\begin{enumerate}
  \item Turunkan \emph{checklist} fungsional 20--30 butir dari GDD kelompok
        Anda, mencakup mekanik, integrasi, dan keseluruhan.
  \item Salin peta level ke \texttt{ch16\_uji\_mandiri.py}, jalankan
        \texttt{python -I ch16\_uji\_mandiri.py --peta namapeta}, lalu perbaiki
        tempat yang dilaporkan terjebak atau hanya dapat menurun.
  \item Bungkus \emph{build} dan ujilah pada komputer lain, bukan komputer
        tempat Anda mengembangkan permainan.
  \item Susun lembar \emph{playtest}; tuliskan bagian yang paling Anda ragukan
        sebagai tujuan sesi.
  \item Laksanakan \emph{playtest} dengan sekurang-kurangnya dua penguji baru;
        jangan memberi petunjuk, dan catat waktu serta tempat berhentinya.
  \item Pindahkan temuan yang dapat diulang ke \emph{log bug}, timbang beratnya,
        perbaiki yang berat, lalu jalankan kembali \emph{checklist} fungsional.
\end{enumerate}
\vspace{0.06cm}
{\footnotesize\color{ubAccent}\itshape Rincian lengkap: bagian praktikum Bab XVI
buku ajar; keluarannya adalah Proyek Akhir. Kumpulkan melalui repositori Git
kelompok.}
\end{frame}

% --------------------------------------------------------- rangkuman
\begin{frame}{Rangkuman}
\begin{itemize}
  \item Pembuat tidak dapat menilai karyanya sendiri; \emph{playtest} menuntut
        penguji yang belum pernah melihat karya, dan Anda tidak boleh
        menolongnya.
  \item Pemeriksaan otomatis dan \emph{playtest} saling melengkapi. Otomatis
        dijalankan lebih dahulu, tetapi tidak dapat menggantikan
        \emph{playtest}.
  \item Pengujian adalah siklus: setiap perbaikan menuntut pemeriksaan ulang dan
        \emph{checklist} yang diperbarui.
  \item \emph{Checklist} fungsional diturunkan dari GDD dan terbagi menjadi
        mekanik, integrasi, serta keseluruhan; kelompok terakhir hanya dapat
        diperiksa pada \emph{build}.
  \item \emph{Soft-lock} tidak menimbulkan pesan kesalahan. Pemeriksaan setiap
        tempat berdiri menemukannya, sedangkan simulasi acak dapat menyesatkan.
  \item Catatan \emph{bug} harus dapat diulang dan ditimbang beratnya; demo
        5--10 menit menampilkan hasil, temuan, dan refleksi teknologi.
\end{itemize}
\end{frame}

% --------------------------------------------------------- cek
\begin{frame}{Cek Pemahaman}
\begin{enumerate}
  \item Mengapa pembuat tidak dapat menilai karyanya sendiri, dan mengapa
        memberi petunjuk saat \emph{playtest} justru menghapus temuan?
  \item Sebuah level memakai $v_0=640$ piksel per detik, $g=1600$ piksel per
        detik kuadrat, dan petak 32 piksel. Hitunglah tinggi lompatan tertinggi
        dalam satuan petak.
  \item Mengapa pemeriksaan setiap tempat berdiri lebih dapat dipercaya daripada
        simulasi 200 titik awal acak? Sebutkan pula dua hal yang tidak dapat
        diperiksa secara otomatis.
\end{enumerate}
\vspace{0.35cm}
{\footnotesize\color{ubMuted}\itshape Jawab tanpa membuka catatan, lalu
bandingkan dengan pasangan di sebelah Anda.}
\end{frame}

% --------------------------------------------------------- penutup
\penutup{Presentasi Proyek Akhir}
  {Minggu 16 adalah UAS: tidak ada materi baru dan tidak ada berkas materi UAS.
   Yang berlangsung adalah uji fungsional dosen atas \emph{build} dan
   \emph{project source} yang telah diserahkan, menggunakan \emph{checklist}
   penilaian, disertai laporan pengembangan dan lembar refleksi pembelajaran.}

\end{document}
